% ----------------------------------------------------------
\chapter{Introdução}
% ----------------------------------------------------------

O fornecimento de água, é um dos serviços mais essenciais a vida humana, seja para abastecer as residências, ou também, serviços públicos, indústrias, usinas hidrelétricas para geração de energia, produção agrícola e agropecuária, e todo tipo de serviço que depende desse recurso natural valioso para a qualidade de vida, a economia e o meio ambiente. Fazer uma gestão eficiente da utilização dos recursos hídricos é primordial para garantir que todos esses setores continuem funcionando de forma sustentável, e é também uma questão de saúde pública, já que segundo informações do ministério da saúde, cenários de seca contribuem para desnutrição por afetar a produção agrícola, desidratação pela baixa disponibilidade de água e transmissão de doenças em cenários que seja necessário acionar fontes de água não devidamente tratadas devido a baixa disponibilidade do sistema de abastecimento.  

De acordo com \cite{ANA_Cantareira}, o Sistema Cantareira (SC), é um dos maiores produtores de água do mundo, se posicionando como o maior entre os sistemas que abastecem a Região Metropolitana de São Paulo (RMSP), é responsável por abastecer cerca de 46\% da população da RMSP. Possui capacidade de volume útil de aproximadamente 980 bilhões de litros e produz uma vazão média de 33m³/s. O sistema é composto pelos reservatórios Jaguari, Jacarei, Cachoeira, Atibainha e Paiva Castro, que são conectados por estruturas de canais e túneis subterrâneos, que formam o chamado Sistema Equivalente Cantareira, que ainda conta com uma estação elevatória para bombeamento de água e a estação de tratamento antes da distribuição e abastecimento. O sistema recebe afluências de diversas bacias hidrográficas, essa vazão afluente é que alimenta os reservatórios com as águas dos mananciais, rios, lençóis freáticos e precipitação. 

A gestão do sistema é de responsabilidade da Agência Nacional de Águas (ANA) e do Departamento de Águas e Energia Elétrica do Estado de São Paulo (DAEE), já a operação do sistema é feita pela Companhia de Saneamento Básico de São Paulo (SABESP), sendo esta a responsável por observar e notificar a necessidade de operações emergenciais em casos excepcionais. Parte das ferramentas de controle do SC, feito pela SABESP, consiste na coleta diária de dados da situação operacional, como: volumes, vazões e precipitação. Estes dados auxiliam no acompanhamento da situação dos reservatórios e da entrada e saída de água no sistema, e através dessas informações, que a SABESP toma decisões sobre ações de racionalização e restrição no abastecimento de água, quando o reservatório atinge níveis alarmantes. 

Analisar o comportamento de varáveis como o volume útil dos reservatórios do Sistema Cantareira, nada mais é, do que de analisar uma série temporal de uma variável, buscando entender o seu comportamento no passado, para projetar cenários futuros. A capacidade de antever a ocorrência de períodos em que o volume útil atinja níveis preocupantes, e tomar ações antecipadas de racionalização no fornecimento de água, pode evitar principalmente que se atinja cenários críticos, tais quais, caracterizem um período de seca ou crise hídrica, como o registrado em 2014 e 2015, onde ocorreu o menor volume armazenado conhecido do SC. Conforme dados da SABESP, em tal crise, o volume útil ficou abaixo de 0\%, o que significa que foi necessário acionar o volume morto do reservatório, cenário extremo, que não se repetiu desde então. 

Diante da importância da gestão de fornecimento de água, nos diversos aspectos e cenários citados acima, o presente trabalho, tem como principal objetivo, criar um sistema capaz de dar suporte a tomada de decisão sobre o abastecimento de água, através da capacidade de visualizar cenários críticos antecipadamente. Para isso, pretende-se construir uma arquitetura que seja capaz de coletar e tratar os dados atualizados pela SABESP diariamente, construir e treinar um modelo de machine learning para aprender sobre o comportamento da variável alvo, o volume útil (\%), e por fim realizar previsões num horizonte determinado, e entregar uma visualização do possível cenário futuro para o comportamento da disponibilidade de água no Sistema Cantareira. 


Daqui pra baixo é o que veio no arquivo de template, deixei pra conseguirmos ver como criar uma figura, uma seção, uma subseção e etc... posteriormente apagamos. 

\section{Recomendações de uso}
Este \textit{template} foi elaborado em \LaTeX. O sumário é gerado automaticamente de acordo com a norma NBR 6027/2012 utilizando a sequência abaixo para diferenciação gráfica nas divisões de seção e subseção.

\vspace{12pt}

\textbf{1 SEÇÃO PRIMÁRIA}

1.1 SEÇÃO SECUNDÁRIA

\textbf{1.1.1 Seção terciária}

\textit{1.1.1.1 Seção quartenária}

1.1.1.1.1 Seção quinária

\vspace{12pt}

\begin{alineas}
    \item Seção primária, use o comando \verb|\section{}|.
    \item Seção secundária, use o comando \verb|\subsection{}|.
    \item Seção terciária, use o comando \verb|\subsubsection{}|.
    \item Seção quartenária, use o comando \verb|\subsubsubsection{}|.
    \item Seção quinária, use o comando \verb|\subsubsubsubsection{}|.
    \item Título das seções de referências, apêndice e anexo são gerados automaticamente pelo \textit{template}.
    \item Para citação com mais de três linhas use o comando \verb|\begin{citacao}|.
    \item Note de rodapé, use o comando \verb|\footnote{}|\footnote{A nota de rodapé é automaticamente formatada pelo \textit{template}.}
\end{alineas}

\section{Objetivos}


Nas seções abaixo estão descritos o objetivo geral e os objetivos específicos deste TCC.


\subsection{Objetivo Geral}


Descrição...


\subsection{Objetivos Específicos}


Descrição...